\exercisesection{\S~9}

All the g-modules considered are assumed to be finite dimensional.

\begin{enumerate}
\def\labelenumi{\arabic{enumi})}
\item
  If \(m\) is an integer \(\geq 0\) , we have \(\dim \mathrm{E}(m\rho) = (m + 1)^{\mathrm{N}}\) , where \(\mathrm{N} = \operatorname{Card}(\mathrm{R}_{+})\) .
\item
  Show that there exists a unique polynomial function \(d\) on \(\mathfrak{h}^*\) such that \(d(\lambda) = \dim \mathrm{E}(\lambda)\) for all \(\lambda \in \mathrm{P}_{++}\) ; its degree is \(\operatorname{Card}(\mathrm{R}_+)\) . We have
\end{enumerate}

\[
d (w \lambda - \rho) = \varepsilon (w) d (\lambda - \rho) \quad \text { if } w \in W, \lambda \in \mathfrak {h} ^ {*}.
\]

In particular, the function \(\lambda \mapsto d(\lambda - \rho)^2\) is invariant under W. Deduce that there exists a unique element \(u\) of the centre of \(\mathrm{U}(\mathfrak{g})\) such that \(\chi_{\lambda}(u) = d(\lambda)^2\) for all \(\lambda \in \mathfrak{h}^*\) (apply Th. 2 of §8, no. 5). When \(\mathfrak{g} = \mathfrak{sl}(2,k)\) , we have \(u = C + 1\) , where \(C\) is the element defined in Exerc. 1 of §1.

\begin{enumerate}
\def\labelenumi{\arabic{enumi})}
\setcounter{enumi}{2}
\tightlist
\item
  Let \(k_{1}, \ldots, k_{l}\) be the characteristic degrees of the algebra of invariants of W (cf.~Chap. V, §5).
\end{enumerate}

\begin{enumerate}
\def\labelenumi{\alph{enumi})}
\item
  Show that, for all \(j \geq 1\) , the number of i such that \(k_{i} > j\) is equal to the number of \(\alpha \in R_{+}\) such that \(\langle \rho, H_{\alpha} \rangle = j\) . (Reduce to the case in which R is irreducible, and use Chap. VI, §4, Exerc. 6 c.) (Cf. §5, Exerc. 5 g).
\item
  Deduce the formula
\end{enumerate}

\[
\prod_ {\alpha \in \mathrm{R} _ {+}} \langle \rho , H _ {\alpha} \rangle = \prod_ {i = 1} ^ {l} (k _ {i} - 1)!.
\]

¶4) Assume that g is simple, and denote by \(\gamma\) the element of \(R_{+}\) such that \(H_{\gamma}\) is the highest root of \(R^{\vee}\) ; write \(H_{\gamma} = \sum_{\alpha \in B} n_{\alpha} H_{\alpha}\) . We have \(\langle \rho, H_{\gamma} \rangle = \sum n_{\alpha} = h - 1\) , where h is the Coxeter number of R (Chap. VI, §1, no. 11, Prop. 31).

\begin{enumerate}
\def\labelenumi{\alph{enumi})}
\tightlist
\item
  Let \(\alpha \in \mathrm{B}\) . Show that, for all \(\beta \in \mathrm{R}_{+}\) ,
\end{enumerate}

\[
\langle \varpi_ {\alpha} + \rho , H _ {\beta} \rangle \leq h + n _ {\alpha} - 1,
\]

and that equality holds if \(\beta = \gamma\) . Deduce that every prime factor of \(\dim \mathrm{E}(\varpi_{\alpha})\) is \(\leq h + n_{\alpha} - 1\) .

\begin{enumerate}
\def\labelenumi{\alph{enumi})}
\setcounter{enumi}{1}
\tightlist
\item
  Assume that \(\varpi_{\alpha}\) is not minuscule, i.e.~\(n_{\alpha} \geq 2\) . Let \(m \in [2, n_{\alpha}]\) and \(p = h + m - 1\) . Verify (cf.~Chap. VI, Plates) that there exists \(\beta \in \mathrm{R}_{+}\) such that \(\langle \varpi_{\alpha}, H_{\beta} \rangle = n_{\alpha}\) and \(\langle \rho, H_{\beta} \rangle = h - 1 - (n_{\alpha} - m)\) , hence \(\langle \varpi_{\alpha} + \rho, H_{\beta} \rangle = p\) . Deduce that, if \(p\) is prime, \(p\) divides \(\dim \mathrm{E}(\varpi_{\alpha})\) . (Remark that \(p\) does not divide any of the \(\langle \rho, H_{\beta} \rangle\) for \(\beta \in \mathrm{R}_{+}\) , cf.~Exerc. 3.)
\end{enumerate}

\(c\) ) When \(\mathfrak{g}\) is of type \(\mathrm{G}_2\) (resp. \(\mathrm{F}_4,\mathrm{E}_8\) ), we have \(h = 6\) (resp. 12, 30), and \(\dim \mathrm{E}(\varpi_{\alpha})\) is divisible by 7 (resp. 13, 31), When \(\mathfrak{g}\) is of type \(\mathrm{E}_6\) (resp. \(\mathrm{E}_7\) ), and \(\varpi_{\alpha}\) is not minuscule, \(\dim \mathrm{E}(\varpi_{\alpha})\) is divisible by 13 (resp. 19).

¶ 5) a) Let \(\alpha \in \mathrm{R}, x \in \mathfrak{g}^{\alpha}, y \in \mathfrak{g}^{-\alpha}\) , and let E be a \(\mathfrak{g}\) -module. Show that, for all \(\lambda \in \mathrm{P}\) ,

\[
\operatorname{Tr} ((x y) _ {\mathrm{E}} | \mathrm{E} ^ {\lambda}) = \operatorname{Tr} ((x y) _ {\mathrm{E}} | \mathrm{E} ^ {\lambda + \alpha}) + \lambda ([ x, y ]) \dim \mathrm{E} ^ {\lambda}.
\]

Deduce that

\[
\sum_ {\lambda \in P} \lambda ([ x, y ]) \dim E ^ {\lambda}. e ^ {\lambda} = (1 - e ^ {- \alpha}) \sum_ {\lambda \in P} \operatorname{Tr} ((x y) _ {E} | E ^ {\lambda}). e ^ {\lambda}.
\]

\begin{enumerate}
\def\labelenumi{\alph{enumi})}
\setcounter{enumi}{1}
\tightlist
\item
  Give \(\mathfrak{h}^*\) a non-degenerate W-invariant symmetric bilinear form \(\langle ,\rangle\) . Let \(\Delta\) be the endomorphism of the vector space \(k[\mathrm{P}]\) such that \(\Delta(e^{\mu}) = \langle \mu ,\mu \rangle e^{\mu}\) for all \(\mu \in \mathrm{P}\) ; if \(a,b\in k[\mathrm{P}]\) , put
\end{enumerate}

\[
\Delta^ {\prime} (a, b) = \Delta (a b) - a \Delta (b) - b \Delta (a).
\]

Prove that

\[
\Delta (\mathrm{J} (e ^ {\mu})) = \langle \mu , \mu \rangle \mathrm{J} (e ^ {\mu}) \quad \text { for } \mu \in \mathrm{P},
\]

\[
\Delta^ {\prime} (e ^ {\lambda}, e ^ {\mu}) = 2 \langle \lambda , \mu \rangle e ^ {\lambda + \mu} \quad \mathrm{for} \lambda , \mu \in \mathrm{P},
\]

\[
\Delta^ {\prime} (a b, c) = a \Delta^ {\prime} (b, c) + b \Delta^ {\prime} (a, c) \quad \text { for } a, b, c \in k [ P ].
\]

\(c)\) Let \(\lambda \in \mathrm{P}_{++}\) , \(c_{\lambda} = \operatorname{ch}(\operatorname{E}(\lambda))\) and \(d = \mathrm{J}(e^{\rho})\) . Prove that

\[
\Delta (c _ {\lambda} d) = \langle \lambda + \rho , \lambda + \rho \rangle c _ {\lambda} d.
\]

(Use \(a\) ), \(b\) ), §6, Cor. of Prop. 7, and Chap. VI, §3, no. 3, formula (3).)

\(d)\) Deduce another proof of the formula of H. Weyl from the above and §7, no. 2, Prop. 5 (iii).

\begin{enumerate}
\def\labelenumi{\alph{enumi})}
\setcounter{enumi}{4}
\tightlist
\item
  For all \(\lambda \in \mathfrak{h}^*\) , put \(\dim \mathrm{E}^{\lambda} = m(\lambda)\) . Deduce from \(a\) that
\end{enumerate}

\[
\begin{array}{c} \operatorname{Tr} ((x y) _ {\mathrm{E}} | \mathrm{E} ^ {\lambda}) = \sum_ {i = 0} ^ {+ \infty} (\lambda + i \alpha) ([ x, y ]) m (\lambda + i \alpha) \\ \sum_ {i = - \infty} ^ {+ \infty} (\lambda + i \alpha) ([ x, y ]) m (\lambda + i \alpha) = 0. \end{array}
\]

\(f)\) Let \(\langle\cdot,\cdot\rangle\) be a non-degenerate invariant symmetric bilinear form on \(\mathfrak{g}\) whose restriction to \(\mathfrak{h}\) is the inverse of the form chosen above. Let \(\Gamma\) be the corresponding Casimir element. Assume that E is simple; put \(\Gamma_{\mathrm{E}}=\gamma.1\) , where \(\gamma\in k\) . By using \(e\) and §2, no. 3, Prop. 6, show that

\[
\gamma m (\lambda) = \langle \lambda , \lambda \rangle m (\lambda) + \sum_ {\alpha \in \mathrm{R}} \sum_ {i = 0} ^ {+ \infty} \langle \lambda + i \alpha , \alpha \rangle m (\lambda + i \alpha)
\]

for all \(\lambda \in \mathfrak{h}^*\) , and then that

\[
\begin{array}{l} \gamma m (\lambda) = \langle \lambda , \lambda \rangle m (\lambda) + \sum_ {\alpha \in \mathrm{R} _ {+}} m (\lambda) \langle \lambda , \alpha \rangle + 2 \sum_ {\alpha \in \mathrm{R} _ {+}} \sum_ {i = 1} ^ {+ \infty} m (\lambda + i \alpha) \langle \lambda + i \alpha , \alpha \rangle \\ = \langle \lambda , \lambda + 2 \rho \rangle m (\lambda) + 2 \sum_ {\alpha \in \mathrm{R} _ {+}} \sum_ {i = 1} ^ {+ \infty} m (\lambda + i \alpha) \langle \lambda + i \alpha , \alpha \rangle . \end{array}
\]

\(g)\) Continue to assume that E is simple; let \(\omega\) be its highest weight. Deduce from \(f)\) that, for all \(\lambda \in \mathfrak{h}^*\) ,

\[
(\langle \omega + \rho , \omega + \rho \rangle - \langle \lambda + \rho , \lambda + \rho \rangle) m (\lambda) = 2 \sum_ {\alpha \in \mathrm{R} _ {+}} \sum_ {i = 1} ^ {+ \infty} m (\lambda + i \alpha) \langle \lambda + i \alpha , \alpha \rangle .
\]

(Recall that, by Prop. 5 of §7, \(\langle\omega+\rho,\omega+\rho\rangle>\langle\lambda+\rho,\lambda+\rho\rangle\) if \(\lambda\) is a weight of E distinct from \(\omega\) . The preceding formula thus gives a procedure for calculating \(m(\lambda)\) step by step.) \(^{18}\)

\begin{enumerate}
\def\labelenumi{\arabic{enumi})}
\setcounter{enumi}{5}
\tightlist
\item
  Let \(x \mapsto x^*\) be the involution of \(k[\mathrm{P}]\) that takes \(e^p\) to \(e^{-p}\) for all \(p \in \mathrm{P}\) . \(a)\) Put \(\mathrm{D} = d^* d = \prod_{\alpha \in \mathrm{R}} (1 - e^\alpha)\) . Show that \(d^* = (-1)^{\mathrm{N}} d\) , where \(\mathrm{N} = \mathrm{Card}(\mathrm{R}_+)\) , and hence that \(\mathrm{D} = (-1)^{\mathrm{N}} d^2\) .
\end{enumerate}

\begin{enumerate}
\def\labelenumi{\alph{enumi})}
\setcounter{enumi}{1}
\tightlist
\item
  Define two linear forms \(\varepsilon\) and I on \(k[\mathrm{P}]\) by the formulas:
\end{enumerate}

\[
\varepsilon (1) = 1, \quad \varepsilon (e ^ {p}) = 0 \quad \text { if } p \in \mathrm{P} - \{0 \}
\]

and

\[
\mathrm{I} (f) = \frac {1}{m} \varepsilon (\mathrm{D}. f), \quad \mathrm{where} m = \mathrm{Card} (\mathrm{W}).
\]

Show, by using the formula \(d = \mathrm{J}(e^{\rho})\) , that \(\mathrm{I}(1) = 1\) .

\begin{enumerate}
\def\labelenumi{\alph{enumi})}
\setcounter{enumi}{2}
\tightlist
\item
  Let \(\lambda \in \mathrm{P}_{++}\) and \(c_{\lambda} = \operatorname{ch} \mathrm{E}(\lambda) = \mathrm{J}(e^{\lambda + \rho}) / d\) . Show that \(\mathrm{I}(c_{\lambda}) = 0\) if \(\lambda \neq 0\) . (Same method as for \(b\) .)
\end{enumerate}

\(d\) ) Show that I takes integer values on the subalgebra \(\mathbf{Z}[\mathrm{P}]^{\mathrm{W}} = \operatorname {ch}\operatorname {R}(\mathfrak{g})\) of \(k[\mathrm{P}]\) . If E is a \(\mathfrak{g}\) -module, the dimension of the space of invariants of \(\mathfrak{g}\) in E is equal to I(ch E). (Reduce to the case in which E is simple and use \(b\) ) and \(c\) .) \(e\) ) We have \(\dim \mathrm{E} = \sum_{\lambda \in \mathrm{P}_{++}}\mathrm{I}(c_\lambda^*\mathrm{ch}\mathrm{E})d(\lambda)\) , where \(d(\lambda) = \dim \mathrm{E}(\lambda)\) . In particular:

\[
\mathrm{I} (c _ {\lambda} ^ {*} c _ {\mu}) = \delta_ {\lambda \mu} \quad \text { if } \lambda , \mu \in \mathrm{P} _ {+ +}.
\]

\(f)\) If \(\lambda, \mu, \nu \in \mathrm{P}_{++}\) , the integer \(m(\lambda, \mu, \nu)\) of Prop. 2 is equal to \(\mathrm{I}(c_{\lambda} c_{\mu} c_{\nu}^{*})\) . Deduce the identity

\[
d (\lambda) d (\mu) = \sum_ {\nu \in \mathrm{P} _ {+ +}} m (\lambda , \mu , \nu) d (\nu) \quad \lambda , \mu , \nu \in \mathrm{P} _ {+ +}.
\]

(Apply e) to the g-module \(\mathrm{E} = \mathrm{E}(\lambda) \otimes \mathrm{E}(\mu).\)

¶ 7) We retain the notations of the proof of Th. 2.

\begin{enumerate}
\def\labelenumi{\alph{enumi})}
\tightlist
\item
  Show that
\end{enumerate}

\[
f _ {\rho} (\mathrm{J} (e ^ {\mu})) = \prod_ {\alpha \in \mathrm{R} _ {+}} (e ^ {(\mu | \alpha) \mathrm{T} / 2} - e ^ {- (\mu | \alpha) \mathrm{T} / 2}).
\]

\begin{enumerate}
\def\labelenumi{\alph{enumi})}
\setcounter{enumi}{1}
\tightlist
\item
  Take \((\cdot |\cdot)\) to be the canonical bilinear form \(\varPhi_{\mathrm{R}}\) (Chap. VI, §1, no. 12). Show that
\end{enumerate}

\[
  f _ {\rho} (\mathrm{J} (e ^ {\mu})) \equiv d _ {\mu} \mathrm{T} ^ {\mathrm{N}} \left(1 + \frac {\mathrm{T} ^ {2}}{48} (\mu | \mu)\right) \pmod {\mathrm{T} ^ {\mathrm{N} + 3} \mathbf {R} [ [ \mathrm{T} ] ]},
\]

where \(d_{\mu} = \prod_{\alpha \in \mathrm{R}_{+}}(\mu |\alpha)\) .

\begin{enumerate}
\def\labelenumi{\alph{enumi})}
\setcounter{enumi}{2}
\tightlist
\item
  From b) and the equality \(\mathrm{J}(e^{\lambda+\rho})=\mathrm{ch}(\mathrm{E}).\mathrm{J}(e^{\rho})\) , deduce the formula
\end{enumerate}

\[
  \sum_ {\mu \in \mathrm{P}} (\mu | \rho) ^ {2} \dim \mathrm{E} ^ {\mu} = \frac {\dim \mathrm{E}}{24} (\lambda | \lambda + 2 \rho).
\]

\(d\) ) Assume that \(\mathfrak{g}\) is simple. Show that \((\rho |\rho) = \dim \mathfrak{g} / 24\) . (Apply \(c\) ) with \(\lambda\) the highest root of R, and use Exerc. 3 of §6.)

¶8) Let \(\psi\) be a polynomial function on \(\mathfrak{h}^*\) of degree \(r\) . Show that there exists a unique polynomial function \(\varPsi\) on \(\mathfrak{h}^*\) that is invariant under W, of degree \(\leq r\) , and such that

\[
\sum_ {\mu \in \mathrm{P}} \psi (\mu) \dim \mathrm{E} ^ {\mu} = \varPsi (\lambda + \rho) \dim \mathrm{E}
\]

for any simple g-module E of highest weight \(\lambda\) .

(Treat first the case in which \(\psi(\mu)=(\mu|\nu)^{r}\) , where \(\nu\in P\) is not orthogonal to any root; for this use the homomorphism \(f_{\nu}\) in the proof of Th. 2 and Chap. V, §5, no. 4, Prop. 5 (i).)

\begin{enumerate}
\def\labelenumi{\arabic{enumi})}
\setcounter{enumi}{8}
\tightlist
\item
  We use the notations of Chap. VI, Plate I, in the case of an algebra \(\mathfrak{g}\) of type \(\mathrm{A}_2\) . Let \(n, p\) be integers \(\geq 0\) .
\end{enumerate}

\[
a) \mathfrak {P} (n \alpha_ {1} + p \alpha_ {2}) = 1 + \inf (n, p).
\]

\begin{enumerate}
\def\labelenumi{\alph{enumi})}
\setcounter{enumi}{1}
\tightlist
\item
  Let \(\lambda = n\varpi_1 + p\varpi_2\) . Then \(\dim \mathrm{E}(\lambda) = \frac{1}{2}(n + 1)(p + 1)(n + p + 2)\) . The multiplicity of the weight 0 of \(\mathrm{E}(\lambda)\) is 0 if \(\lambda\) is not radical, i.e.~if \(n \not\equiv p \pmod{3}\) ; if \(\lambda\) is radical, it is \(1 + \inf(n,p)\) .
\end{enumerate}

\begin{enumerate}
\def\labelenumi{\arabic{enumi})}
\setcounter{enumi}{9}
\tightlist
\item
  We use the notations of Chap. VI, Plate II, in the case of an algebra of type \(B_{2}\) . Let n, p be integers \(\geq 0\) .
\end{enumerate}

\begin{enumerate}
\def\labelenumi{\alph{enumi})}
\tightlist
\item
  We have
\end{enumerate}

\[
\begin{array}{c} \mathfrak {P} (n \alpha_ {1} + p (\alpha_ {1} + 2 \alpha_ {2})) = 1 + \frac {1}{2} p (p + 3) \\ \mathfrak {P} (n \alpha_ {2} + p (\alpha_ {1} + \alpha_ {2})) = [ p ^ {2} / 4 ] + p + 1 \\ \mathfrak {P} (n (\alpha_ {1} + \alpha_ {2}) + p (\alpha_ {1} + 2 \alpha_ {2})) = [ n ^ {2} / 4 ] + n + 1 + n p + \frac {1}{2} p (p + 3). \end{array}
\]

\begin{enumerate}
\def\labelenumi{\alph{enumi})}
\setcounter{enumi}{1}
\tightlist
\item
  Let \(\lambda = n(\alpha_{1} + \alpha_{2}) + p(\alpha_{1} + 2\alpha_{2})\) . The multiplicity of the weight 0 in \(\operatorname{E}(\lambda)\) is
\end{enumerate}

\[
[ n / 2 ] + 1 + n p + p.
\]

\begin{enumerate}
\def\labelenumi{\arabic{enumi})}
\setcounter{enumi}{10}
\item
  Assume that \(\mathfrak{g}\) is not a product of algebras of rank 1. Let \(n\in \mathbf{N}\) . Show that there exists a simple \(\mathfrak{g}\) -module one of whose weights is of multiplicity \(\geq n\) . (Suppose not. Let \(E_{\lambda}\) be the simple g-module of highest weight \(\lambda \in P_{++}\) . Compare \(\dim E_{\lambda}\) and the number of distinct weights of \(E_{\lambda}\) when \((\lambda|\lambda) \to \infty\) (with the notations of Th. 2).)
\item
  Let \(U^{0}\) be the commutant of h in \(\mathrm{U}(\mathfrak{g})\) . If g is of rank 1, \(U^{0}\) is commutative. If g is of rank \(\geq 2\) , \(U^{0}\) admits simple representations of arbitrarily large finite dimension. (Use Exerc. 11 and Exerc. 23 a) of §7.)
\item
  Let R be a root system in a vector space V. Two elements \(v_{1}, v_{2}\) of V are said to be disjoint if R is the direct sum of two root systems \(R_{1}\) and \(R_{2}\) (Chap. VI, §1, no. 2) such that \(v_{i}\) belongs to the vector subspace of V generated by \(R_{i}, i = 1, 2\) . Show that two elements of \(V_{R}\) that belong to the same chamber of R are disjoint if and only if they are orthogonal.
\end{enumerate}

¶ 14) a) Let \(\mu, \nu \in \mathrm{P}_{++}\) and let \(\gamma\) be a weight of \(\mathrm{E}(\mu)\) . Let \(\rho_{\mu}\) be the representation of \(\mathfrak{g}\) on \(\mathrm{E}(\mu)\) . Let \(X_{\alpha} \in \mathfrak{g}^{\alpha} - \{0\}, Y_{\alpha} \in \mathfrak{g}^{-\alpha} - \{0\}\) . If \(\alpha \in \mathrm{B}\) , put \(\nu_{\alpha} = \nu(H_{\alpha})\) , and

\[
\mathrm{E} ^ {+} (\mu , \gamma , \nu) = \mathrm{E} (\mu) ^ {\gamma} \cap \bigcap_ {\alpha \in \mathrm{B}} \mathrm{Ker} \rho_ {\mu} (X _ {\alpha}) ^ {\nu_ {\alpha} + 1},
\]

\[
\mathrm{E} ^ {-} (\mu , \gamma , \nu) = \mathrm{E} (\mu) ^ {\gamma} \cap \bigcap_ {\alpha \in \mathrm{B}} \operatorname{Ker} \rho_ {\mu} (Y _ {\alpha}) ^ {\nu_ {\alpha} + 1},
\]
% semantic-fragments: legacy-input-seam
\relax{}
\[
d ^ {+} (\mu , \gamma , \nu) = \dim \mathrm{E} ^ {+} (\mu , \gamma , \nu), \quad d ^ {-} (\mu , \gamma , \nu) = \dim \mathrm{E} ^ {-} (\mu , \gamma , \nu).
\]

For all \(\lambda \in \mathfrak{h}^*\) , put \(\lambda^{*} = -w_{0}\lambda\) , where \(w_{0}\) is the element of W that takes B to \(-B\) . Show that

\[
d ^ {+} (\mu , \gamma , \nu) = d ^ {-} (\mu , - \gamma^ {*}, \nu^ {*}).
\]

\begin{enumerate}
\def\labelenumi{\alph{enumi})}
\setcounter{enumi}{1}
\item
  Let \(\lambda_1, \lambda_2 \in \mathrm{P}_{++}\) , V the \(\mathfrak{g}\) -module \(\mathrm{Hom}_k(\mathrm{E}(\lambda_1^*), \mathrm{E}(\lambda_2))\) , U the set of \(\varphi \in V\) such that \(Y_\alpha. \varphi = 0\) for all \(\alpha \in B\) , and \(w\) a primitive vector in \(\mathrm{E}(\lambda_1^*)\) . Show that \(\varphi \mapsto \varphi(w)\) is an isomorphism from U to the set of \(v \in \mathrm{E}(\lambda_2)\) such that \(Y_\alpha^{\lambda_1^*(H_\alpha)+1}. v = 0\) for all \(\alpha \in B\) . (Use Exerc. 15 of §7 to prove surjectivity.)
\item
  With the notations of Prop. 2, prove that, for \(\lambda_1, \lambda_2, \lambda \in \mathrm{P}_{++}\) ,
\end{enumerate}

\[
\begin{array}{r} m (\lambda_ {1}, \lambda_ {2}, \lambda) = d ^ {+} (\lambda , \lambda_ {2} - \lambda_ {1} ^ {*}, \lambda_ {1} ^ {*}) = d ^ {-} (\lambda , \lambda_ {1} - \lambda_ {2} ^ {*}, \lambda_ {1}) \\ = d ^ {+} (\lambda_ {1}, \lambda - \lambda_ {2}, \lambda_ {2}) = d ^ {-} (\lambda_ {1}, \lambda_ {2} ^ {*} - \lambda^ {*}, \lambda_ {2} ^ {*}). \end{array}
\]

(Observe that \(m(\lambda_1,\lambda_2,\lambda)\) is the dimension of the space of \(\mathfrak{g}\) -invariant elements of

\[
\mathrm{E} (\lambda) ^ {*} \otimes \mathrm{E} (\lambda_ {1}) \otimes \mathrm{E} (\lambda_ {2}),
\]

and hence is equal to \(m(\lambda_1^*, \lambda, \lambda_2)\) .)

\(d)\) Let \(\lambda_1, \lambda_2 \in \mathrm{P}_{++}\) , and let \(\lambda\) be the unique element of \(\mathrm{P}_{++} \cap \mathrm{W}.\) ( \(\lambda_1 - \lambda_2^*\) ). We have

\[
m (\lambda_ {1}, \lambda_ {2}, \lambda) = 1.
\]

(Use c.) Deduce that \(\mathrm{E}(\lambda_1) \otimes \mathrm{E}(\lambda_2)\) contains a unique submodule isomorphic to \(\mathrm{E}(\lambda)\) .

\begin{enumerate}
\def\labelenumi{\alph{enumi})}
\setcounter{enumi}{4}
\tightlist
\item
  We retain the notations of \(d\) ). Show the equivalence of the following conditions:
\end{enumerate}

\begin{enumerate}
\def\labelenumi{(\roman{enumi})}
\item
  \(\lambda = \lambda_{1} + \lambda_{2}\)
\item
  \(\| \lambda \| = \| \lambda_1 + \lambda_2 \|\)
\item
  \(\lambda_{1}\) and \(\lambda_2^*\) are orthogonal
\item
  \(\lambda_{1}\) and \(\lambda_2^*\) are disjoint (Exerc. 13)
\item
  \(\lambda_{1}\) and \(\lambda_{2}\) are disjoint.
\end{enumerate}

Conclude that \(\mathrm{E}(\lambda_{1})\otimes\mathrm{E}(\lambda_{2})\) is not a simple module unless \(\lambda_{1}\) and \(\lambda_{2}\) are disjoint (and hence another proof of Exerc. 18 f) of §7).

\begin{enumerate}
\def\labelenumi{\arabic{enumi})}
\setcounter{enumi}{14}
\tightlist
\item
  Put \(\mathrm{N} = \mathrm{Card}(\mathrm{R}_{+})\) , \(c = \prod_{\alpha \in \mathrm{R}_{+}} \langle \rho, H_{\alpha} \rangle\) , and \(d(\lambda) = \dim \mathrm{E}(\lambda)\) if \(\lambda \in \mathrm{P}_{++}\) . Show that, for any real number \(s > 0\) ,
\end{enumerate}

\[
\sum_ {\lambda \in \mathrm{P} _ {+ +}} d (\lambda) ^ {- s} \leq \frac {1}{c} \left(\sum_ {m = 1} ^ {\infty} m ^ {- s}\right) ^ {\mathrm{N}}.
\]

Deduce that \(\sum_{\lambda \in \mathrm{P}_{+ + }}d(\lambda)^{-s} <   + \infty\) if \(s > 1\)

¶ 16) a) Take g to be of type \(F_{4}\) and use the notations of Chap. VI, Plate VIII. If i = 1, 2, 3, 4, put

\[
\mathcal {X} _ {i} = \mathrm{P} _ {+ +} \cap (\varpi_ {i} - \mathrm{Q} _ {+}).
\]

The set of weights of \(\mathrm{E}(\varpi_i)\) is the disjoint union of the \(\mathrm{W}\omega\) , where \(\omega\) belongs to \(\mathcal{X}_i\) (cf.~§7, Prop. 5 (iv)). We have:

\[
\mathcal {X} _ {1} = \{0, \varpi_ {1}, \varpi_ {4} \};
\]

\[
\mathcal {X} _ {2} = \{0, \varpi_ {1}, \varpi_ {2}, \varpi_ {3}, \varpi_ {4}, \varpi_ {1} + \varpi_ {4}, 2 \varpi_ {4} \};
\]

\[
\mathcal {X} _ {3} = \{0, \varpi_ {1}, \varpi_ {3}, \varpi_ {4} \};
\]

\[
\mathcal {X} _ {4} = \{0, \varpi_ {4} \}.
\]

\begin{enumerate}
\def\labelenumi{\alph{enumi})}
\setcounter{enumi}{1}
\tightlist
\item
  Show, by means of Th. 2, that:
\end{enumerate}

\(\dim \mathrm{E}(\varpi_1) = 52,\quad \dim \mathrm{E}(\varpi_2) = 1274,\quad \dim \mathrm{E}(\varpi_3) = 273,\quad \dim \mathrm{E}(\varpi_4) = 26.\)

\begin{enumerate}
\def\labelenumi{\alph{enumi})}
\setcounter{enumi}{2}
\tightlist
\item
  By using Chap. V, §3, Prop. 1, and the plates of Chap. VI, show that
\end{enumerate}

\[
\mathrm{Card} (\mathrm{W} \varpi_ {1}) = 2 ^ {7} 3 ^ {2} 2 ^ {- 3} (3!) ^ {- 1} = 24.
\]

Calculate \(\mathrm{Card}(\mathrm{W}\varpi_{2}),\ldots,\mathrm{Card}(\mathrm{W}.2\varpi_{4})\) similarly. Deduce that the number of weights of \(\mathrm{E}(\varpi_{2})\) is 553; since this number is strictly less than \(\dim\mathrm{E}(\varpi_{2})\) , one of these weights is of multiplicity \(\geq2\) .

\begin{enumerate}
\def\labelenumi{\alph{enumi})}
\setcounter{enumi}{3}
\item
  Make analogous calculations for \(\varpi_{1},\varpi_{3},\varpi_{4}\) . By using Exerc. 21 of §7, deduce that, if \(\rho\) is a non-zero simple representation of g, \(\rho\) admits a weight of multiplicity \(\geq2\) .
\item
  Prove the same result for a simple algebra of type \(\mathrm{E}_8\) .
\end{enumerate}

\(f)\) Let \(\mathfrak{a}\) be a splittable simple Lie algebra. Deduce from \(d\) , \(e)\) and Prop. 7 and 8 of §7 the equivalence of the following properties:

\begin{enumerate}
\def\labelenumi{(\roman{enumi})}
\item
  \(\mathfrak{a}\) admits a non-zero simple representation all of whose weights are of multiplicity 1;
\item
  \(\mathfrak{a}\) is neither of type \(\mathrm{F}_4\) nor of type \(\mathrm{E}_8\) .
\end{enumerate}

